\chapter{Potok Ricciego przed osobliwościami}
\label{ch:potok-ricciego}
\index{potok Ricciego}

Metryka pozwala mierzyć długości i objętości, a tensor Ricciego
opisuje część jej krzywizny. Można więc zapytać, co się stanie,
gdy metrykę będziemy zmieniać z prędkością wyznaczoną przez
ten tensor. Na sferze oczekujemy kurczenia, ponieważ jej Ricci
jest dodatni; płaski torus powinien pozostać nieruchomy.
Oba przewidywania sprawdzimy rachunkiem. Wyprowadzimy równania
dla miary i krzywizny, porównamy potok zwykły ze znormalizowanym
i obliczymy przykłady metryk iloczynowych oraz jednorodnych.
Na końcu ustalimy, co musi się stać z krzywizną, jeżeli gładkiego
rozwiązania nie można przedłużyć przez skończoną chwilę.
Wskażemy dokładnie miejsca, w których potrzebna jest analiza
równań parabolicznych.

W całym rozdziale $M^n$ jest gładką, zwartą rozmaitością bez
brzegu, $n\geq2$, a metryki są riemannowskie. Orientacja nie jest
potrzebna do określenia gęstości objętości. Zachowujemy konwencję
krzywizny z Definicji~\ref{def:tensor-krzywizny}: okrągła sfera
ma dodatni tensor Ricciego. Piszemy
\[
 \Delta_g f=\operatorname{div}_g(\nabla^g f)
            =g^{ij}\nabla_i\nabla_j f.
\]
Jest to konwencja z rozdziału~\ref{ch:krzywizna}; na zwartej
rozmaitości $\int_M f\Delta_g f\,\dd\mu_g
=-\int_M|\nabla f|_g^2\,\dd\mu_g$. Na orientowalnej $M$
laplasjan Hodge'a z Rozdziału~\ref{ch:hodge-bochner},
$\Delta_H=\dd\delta+\delta\dd$, spełnia na funkcjach
$\Delta_H=-\Delta_g$. To rozróżnienie ustala
znak wszystkich poniższych równań cieplnych.

\section{Równanie i metryki o stałej krzywiźnie}
\label{sec:ricci-flow-definition}

\begin{definicja}[Potok Ricciego Hamiltona]
\label{def:ricci-flow}
Niech $g_0$ będzie gładką metryką na $M$. \emph{Potokiem
Ricciego} o wartości początkowej $g_0$ na przedziale $[0,T)$
nazywamy gładką rodzinę metryk $g(t)$ spełniającą
\begin{equation}\label{eq:ricci-flow}
 \partial_tg(t)=-2\operatorname{Ric}(g(t)),
 \qquad g(0)=g_0.
\end{equation}
Jest to wersja \emph{nienormalizowana}: nie dodajemy składnika,
który utrzymywałby całkowitą objętość na stałym poziomie.
\end{definicja}


\begin{figure}[htbp]
\centering
\begin{tikzpicture}[>=Stealth, font=\small]

% Podpisy etapów
\node[align=center] at (2.20,6.55)
  {\textbf{Metryka początkowa}\\[2pt]$(\mathbb R^3,g(t_0))$};

\node[align=center] at (7.35,6.55)
  {\textbf{Etap pośredni}\\[2pt]$(\mathbb R^3,g(t_1))$};

\node[align=center] at (12.50,6.55)
  {\textbf{Geometria bardziej jednorodna}\\[2pt]$(\mathbb R^3,g(t_2))$};

% Trzy schematyczne powierzchnie siatkowe
% A: wielkość wybrzuszenia
% B: głębokość zwężeń
\foreach \X/\A/\B in {
  0/0.72/0.36,
  5.15/0.34/0.16,
  10.30/0.00/0.00
}{
  \begin{axis}[
    at={(\X cm,1.12cm)},
    anchor=south west,
    width=4.40cm,
    height=4.85cm,
    scale only axis,
    hide axis,
    view={50}{18},
    xmin=-2.05,xmax=2.05,
    ymin=-2.05,ymax=2.05,
    zmin=-1.8,zmax=1.8,
    z buffer=sort,
    samples=34,
    samples y=17,
    domain=0:360,
    y domain=-1.7:1.7,
    colormap={mojniebieski}{
      rgb255=(213,239,253)
      rgb255=(154,205,239)
    }
  ]

  \addplot3[
    surf,
    shader=faceted,
    faceted color=blue!50!black,
    line width=0.16pt,
    opacity=0.78
  ]
  (
    {
      (1+\A*exp(-3*y*y)
      -\B*(exp(-6*(y-1.03)^2)
      +exp(-6*(y+1.03)^2)))*cos(x)
    },
    {
      (1+\A*exp(-3*y*y)
      -\B*(exp(-6*(y-1.03)^2)
      +exp(-6*(y+1.03)^2)))*sin(x)
    },
    {y}
  );

  \end{axis}
}

% Strzałki ewolucji
\draw[->,very thick]
  (4.52,3.47)--(5.02,3.47);

\draw[->,very thick]
  (9.67,3.47)--(10.17,3.47);

% Oś czasu i równanie potoku
\draw[->,thick]
  (0.25,-0.10)--(14.75,-0.10);

\node[fill=white,inner sep=2pt]
  at (7.50,-0.10)
  {
    $\displaystyle
    \partial_t g
    =-2\operatorname{Ric}(g),
    \qquad t_0<t_1<t_2
    $
  };

\end{tikzpicture}

\caption{Schemat ideowy zmiany geometrii na
$\mathbb R^3$ w potoku Ricciego
$\partial_t g=-2\operatorname{Ric}(g)$.
Niebieskie powierzchnie siatkowe symbolizują
zmniejszanie się wybranej lokalnej deformacji;
nie są reprezentacją całej trójwymiarowej metryki.
Rysunek nie zakłada, że dowolny potok Ricciego
zbiega do metryki euklidesowej.}
\label{fig:ricci-r3-schemat}
\end{figure}


Dla wektora $v\in T_pM$, trzymanego nieruchomo przy zmianie
czasu, równanie daje
\[
 \frac{\dd}{\dd t}|v|_{g(t)}^2
 =-2\operatorname{Ric}_{g(t)}(v,v).
\]
Dodatni Ricci zmniejsza więc kwadrat długości w danym
kierunku; ujemny go zwiększa. Nazwa „potok” oznacza tu
ewolucję tensora metrycznego, a nie przepływ punktów
rozmaitości z rozdziału~\ref{ch:pola}.

\begin{stwierdzenie}[Stała zmiana skali]
\label{prop:ricci-scaling}
Jeżeli $g_0$ jest metryką i $c>0$ jest liczbą niezależną
od punktu $M$, to dla $\widetilde g=cg_0$ zachodzi
\[
 \nabla^{\widetilde g}=\nabla^{g_0},\qquad
 \operatorname{Ric}(\widetilde g)=\operatorname{Ric}(g_0),
 \qquad \operatorname{Scal}(\widetilde g)
          =c^{-1}\operatorname{Scal}(g_0),
\]
\[
 K_{\widetilde g}=c^{-1}K_{g_0},\qquad
 \dd\mu_{\widetilde g}=c^{n/2}\dd\mu_{g_0}.
\]
Wzór na $K$ dotyczy każdej płaszczyzny stycznej.
\end{stwierdzenie}
\begin{proof}
W mapie $\widetilde g_{ij}=cg_{0,ij}$ i
$\widetilde g^{ij}=c^{-1}g_0^{ij}$. We wzorze na symbole
Christoffela~\eqref{eq:christoffel-metryka} czynnik $c$
z pochodnych metryki znosi $c^{-1}$ z macierzy
odwrotnej. Koneksja i operator $R(X,Y)Z$ są więc
niezmienione. Ricci z rozdziału~\ref{ch:krzywizna}
jest śladem endomorfizmu $W\mapsto R(W,X)Y$, zatem
również się nie zmienia jako tensor typu $(0,2)$.
Skalar wymaga jeszcze kontrakcji z
$\widetilde g^{ij}=c^{-1}g_0^{ij}$.

W ilorazie definiującym krzywiznę sekcyjną
\eqref{eq:krzywizna-sekcyjna} licznik jest liniowy
w metryce i mnoży się przez $c$, a wyznacznik
macierzy Grama w mianowniku przez $c^2$. Wreszcie
$\det(cg_{0,ij})=c^n\det(g_{0,ij})$; pierwiastek
wyznacznika daje ostatnią równość.
\end{proof}

\begin{stwierdzenie}[Metryka Einsteina jako rozwiązanie]
\label{prop:ricci-einstein-flow}
Załóżmy, że $\operatorname{Ric}(g_0)=\lambda g_0$
z pewną stałą $\lambda\in\R$. Wtedy
\begin{equation}\label{eq:ricci-einstein-solution}
 g(t)=(1-2\lambda t)g_0
\end{equation}
spełnia~\eqref{eq:ricci-flow} na każdym przedziale,
na którym $1-2\lambda t>0$. Gdy $\lambda>0$,
rozwiązanie tego kształtu kończy się przy
$t=(2\lambda)^{-1}$; gdy $\lambda=0$, jest stacjonarne;
gdy $\lambda<0$, skala rośnie.
\end{stwierdzenie}
\begin{proof}
Wstawmy $g(t)=c(t)g_0$ z $c(0)=1$. Ponieważ
$c(t)$ jest stałe na $M$, Stwierdzenie~\ref{prop:ricci-scaling}
daje $\operatorname{Ric}(g(t))=\lambda g_0$.
Równanie~\eqref{eq:ricci-flow} przyjmuje postać
$c'(t)g_0=-2\lambda g_0$, czyli $c'=-2\lambda$.
Całkując z warunkiem $c(0)=1$, dostajemy podany wzór.
Dodatniość $c(t)$ jest konieczna, by tensor był metryką.
\end{proof}

\begin{przyklad}[Okrągła sfera]
\label{prz:ricci-sphere}
Niech $g_{\mathrm{jedn}}$ będzie metryką sfery jednostkowej
$S^n$, a $g_0=a_0^2g_{\mathrm{jedn}}$, $a_0>0$.
Z~\eqref{eq:stala-krzywizna-tensor} i
Stwierdzenia~\ref{prop:ricci-scaling} otrzymujemy
\[
 \operatorname{Ric}(g_0)
 =(n-1)g_{\mathrm{jedn}}
 =\frac{n-1}{a_0^2}g_0.
\]
Zatem rozwiązanie z poprzedniego stwierdzenia ma postać
\begin{equation}\label{eq:ricci-sphere}
 g(t)=a(t)^2g_{\mathrm{jedn}},\qquad
 a(t)^2=a_0^2-2(n-1)t,
 \qquad 0\leq t<\frac{a_0^2}{2(n-1)}.
\end{equation}
Promień $a(t)$ maleje, a $K(t)=a(t)^{-2}$ rośnie.
W chwili granicznej tensor traci dodatnią określoność;
nie jest to gładkie przedłużenie potoku przez tę chwilę.
Dla $n=2$ pole powierzchni wynosi
$4\pi a(t)^2=4\pi a_0^2-8\pi t$.
\end{przyklad}

\begin{przyklad}[Płaski torus]
\label{prz:ricci-torus}
Niech $T^n=\R^n/\Lambda$ dla kraty $\Lambda$,
a $g_0$ niech będzie metryką pochodzącą od stałego
dodatnio określonego iloczynu skalarnego na $\R^n$.
W tych współrzędnych symbole Christoffela i tensor
krzywizny są zerowe, więc $\operatorname{Ric}(g_0)=0$.
Rodzina $g(t)=g_0$ spełnia~\eqref{eq:ricci-flow}.
W szczególności płaski torus dwuwymiarowy zachowuje
swój kształt i pole. Sam warunek $\chi(T^2)=0$ nie mówi,
że \emph{każda} metryka torusa jest płaska; dla innych
metryk pole pozostanie stałe, choć metryka może się
zmieniać.
\end{przyklad}

\section{Ewolucja miary i krzywizny skalarnej}
\label{sec:ricci-evolution}

\begin{stwierdzenie}[Zmiana elementu objętości]
\label{prop:ricci-volume}
Dla dowolnej gładkiej rodziny metryk o
$h=\partial_tg$ mamy
\[
 \partial_t\dd\mu_g
   =\tfrac12\operatorname{tr}_g h\,\dd\mu_g.
\]
Wzdłuż potoku Ricciego wynika stąd
\begin{equation}\label{eq:ricci-volume}
 \partial_t\dd\mu_g=-\operatorname{Scal}_g\,\dd\mu_g,
 \qquad
 \frac{\dd}{\dd t}\operatorname{Vol}_{g(t)}(M)
 =-\int_M\operatorname{Scal}_{g(t)}\,\dd\mu_{g(t)}.
\end{equation}
\end{stwierdzenie}
\begin{proof}
W mapie gęstość objętości to
$\sqrt{\det(g_{ij})}\,|\dd x^1\cdots\dd x^n|$.
Wzór Jacobiego na pochodną wyznacznika daje
\[
 \partial_t\det(g_{ij})
 =\det(g_{ij})\operatorname{tr}(g^{-1}h)
 =\det(g_{ij})g^{ij}h_{ij}.
\]
Pochodna pierwiastka wprowadza czynnik $1/2$.
Dla $h=-2\operatorname{Ric}$ ślad wynosi
$-2\operatorname{Scal}$. Zwarta rozmaitość pozwala
różniczkować całkę po $M$ pod znakiem całkowania.
\end{proof}

Do równania dla skalara potrzebujemy jednej tożsamości
wynikającej z krzywizny koneksji Leviego-Civity.

\begin{lemat}[Ściągnięta druga tożsamość Bianchiego]
\label{lem:ricci-contracted-bianchi}
Dla każdej gładkiej metryki riemannowskiej
\begin{equation}\label{eq:ricci-contracted-bianchi}
 \nabla^i\operatorname{Ric}_{ij}
       =\tfrac12\nabla_j\operatorname{Scal}.
\end{equation}
\end{lemat}
\begin{proof}
Użyjmy oznaczeń z~\eqref{eq:riemann-wspolrzedne}:
$R(\partial_i,\partial_j)\partial_b
=R^a{}_{b i j}\partial_a$ i
$\operatorname{Ric}_{bj}=R^a{}_{b a j}$; w drugim
wzorze użyliśmy symetrii Ricciego z
Rozdziału~\ref{ch:krzywizna}. Pokażmy krok prowadzący
do drugiej tożsamości Bianchiego. W punkcie wybieramy
mapę normalną i piszemy $D_i=\nabla_{\partial_i}$.
Ponieważ pola współrzędnych komutują,
$[D_i,D_j]V=R(\partial_i,\partial_j)V$.
Tożsamość Jacobiego dla operatorów $D_i$ ma postać
$[D_k,[D_i,D_j]]+[D_i,[D_j,D_k]]+[D_j,[D_k,D_i]]=0$.
Komutator $D_k$ z działaniem tensora $R_{ij}$
na $V$ jest $\nabla_kR_{ij}$, więc otrzymujemy
\emph{drugą tożsamość Bianchiego}
\[
 \nabla_kR^a{}_{b i j}
 +\nabla_iR^a{}_{b j k}
 +\nabla_jR^a{}_{b k i}=0.
\]
Równość jest tensorowa, więc obowiązuje w każdej
mapie. Kontrakcja $k=a$ daje
\[
 \nabla_aR^a{}_{b i j}
 -\nabla_i\operatorname{Ric}_{bj}
 +\nabla_j\operatorname{Ric}_{bi}=0.
\]
Po kontrakcji $b=i$ pierwszy składnik jest
$-\nabla^a\operatorname{Ric}_{aj}$, ponieważ symetrie
Riemanna dają $g^{bi}R^a{}_{b i j}
=-\operatorname{Ric}^a{}_j$.
Pozostałe składniki to odpowiednio
$-\nabla^i\operatorname{Ric}_{ij}$ i
$\nabla_j\operatorname{Scal}$. Ich suma jest zerem,
co daje~\eqref{eq:ricci-contracted-bianchi}.
\end{proof}

\begin{twierdzenie}[Ewolucja krzywizny skalarnej]
\label{thm:ricci-scalar-evolution}
Jeżeli $g(t)$ jest gładkim potokiem Ricciego, to
\begin{equation}\label{eq:ricci-scalar-evolution}
 \partial_t\operatorname{Scal}
 =\Delta_g\operatorname{Scal}
  +2|\operatorname{Ric}|_g^2
 =-\Delta_H\operatorname{Scal}
  +2|\operatorname{Ric}|_g^2.
\end{equation}
Drugą równość czytamy tylko w stopniu zero, gdzie
$\Delta_H=-\Delta_g$.
\end{twierdzenie}
\begin{proof}
Najpierw obliczmy zmianę skalara przy dowolnej
rodzinie $g(t)$ i oznaczmy $h_{ij}=\partial_tg_{ij}$.
Różniczkując $g^{ik}g_{kj}=\delta^i_j$, dostajemy
\begin{equation}\label{eq:ricci-inverse-variation}
 \partial_tg^{ij}=-g^{ip}h_{pq}g^{qj}=-h^{ij}.
\end{equation}
Różnica dwóch koneksji jest tensorem. Zróżniczkowanie
wzoru na koneksję Leviego-Civity i użycie $\nabla g=0$
dają
\begin{equation}\label{eq:ricci-connection-variation}
 \partial_t\Gamma^k_{ij}
 =\tfrac12g^{k\ell}
   (\nabla_i h_{j\ell}+\nabla_j h_{i\ell}
                        -\nabla_\ell h_{ij}).
\end{equation}
Z definicji Ricciego przez krzywiznę pochodna
składników kwadratowych w $\Gamma$ łączy się
z pochodnymi zwykłymi w pochodne kowariantne:
\[
 \partial_t\operatorname{Ric}_{ij}
 =\nabla_k(\partial_t\Gamma^k_{ij})
   -\nabla_i(\partial_t\Gamma^k_{kj}).
\]
Ponieważ $\partial_t\Gamma^k_{kj}
=\tfrac12\nabla_j(g^{pq}h_{pq})$, po kontrakcji
z $g^{ij}$ pierwszego składnika otrzymujemy
$\nabla_k\nabla_i h^{ik}
-\tfrac12\Delta_g(\operatorname{tr}_gh)$,
a z drugiego
$-\tfrac12\Delta_g(\operatorname{tr}_gh)$.
Różniczkując dodatkowo kontrakcję
$\operatorname{Scal}=g^{ij}\operatorname{Ric}_{ij}$
i korzystając z~\eqref{eq:ricci-inverse-variation},
otrzymujemy pełny wzór wariacyjny
\begin{equation}\label{eq:ricci-scalar-variation}
 \partial_t\operatorname{Scal}
 =-\langle h,\operatorname{Ric}\rangle_g
   +\nabla_k\nabla_i h^{ik}
   -\Delta_g(\operatorname{tr}_g h).
\end{equation}

Teraz $h=-2\operatorname{Ric}$, więc pierwszy składnik
jest $2|\operatorname{Ric}|^2$, a
$\operatorname{tr}_gh=-2\operatorname{Scal}$.
Lemat~\ref{lem:ricci-contracted-bianchi} daje
\[
 \nabla_k\nabla_i h^{ik}
 =-2\nabla_k\nabla_i\operatorname{Ric}^{ik}
 =-\Delta_g\operatorname{Scal}.
\]
Ostatni składnik~\eqref{eq:ricci-scalar-variation}
jest $+2\Delta_g\operatorname{Scal}$; suma obu
wyrazów laplasjanowych to $+\Delta_g\operatorname{Scal}$.
\end{proof}

\section{Powierzchnie, pole i zasada maksimum}
\label{sec:ricci-surfaces}
\index{zasada maksimum!dla krzywizny skalarnej}

W wymiarze dwa mamy już
$\operatorname{Ric}=Kg$ i $\operatorname{Scal}=2K$
ze Stwierdzenia w Sekcji~\ref{sec:ricci-krzywizna-skalarna}.
Równanie potoku przyjmuje wtedy postać
\begin{equation}\label{eq:ricci-surface-metric}
 \partial_tg=-2Kg=-\operatorname{Scal}\,g.
\end{equation}
W danej chwili zmiana metryki jest więc proporcjonalna
do samej metryki: zachowuje lokalne kąty, choć współczynnik
skali zależy od punktu. Ponadto
$|\operatorname{Ric}|^2=|Kg|^2=2K^2$, ponieważ
$g^{ia}g^{jb}g_{ij}g_{ab}=2$. Dzieląc
\eqref{eq:ricci-scalar-evolution} przez dwa, dostajemy
\begin{equation}\label{eq:ricci-surface-curvature}
 \partial_tK=\Delta_gK+2K^2.
\end{equation}

\begin{stwierdzenie}[Pole zamkniętej powierzchni]
\label{prop:ricci-surface-area}
Niech $M^2$ będzie zwartą powierzchnią bez brzegu,
a $g(t)$ gładkim potokiem Ricciego. Wtedy
\begin{equation}\label{eq:ricci-surface-area}
 \operatorname{Area}_{g(t)}(M)
 =\operatorname{Area}_{g_0}(M)-4\pi\chi(M)t
\end{equation}
na całym przedziale istnienia rozwiązania.
\end{stwierdzenie}
\begin{proof}
Ze Stwierdzenia~\ref{prop:ricci-volume} i równości
$\operatorname{Scal}=2K$ wynika
\[
 \frac{\dd}{\dd t}\operatorname{Area}_{g(t)}(M)
 =-2\int_M K_{g(t)}\,\dd A_{g(t)}.
\]
Dopiero teraz stosujemy udowodzone dla powierzchni
zamkniętej twierdzenie Gaussa--Bonneta,
Wniosek~\ref{cor:gauss-bonnet-zamknieta}:
$\int_MK\,\dd A=2\pi\chi(M)$.
Zatem pochodna pola jest stałą $-4\pi\chi(M)$;
całkowanie po czasie daje~\eqref{eq:ricci-surface-area}.
Ten sam związek z charakterystyką Eulera wynika z
dwuwymiarowego przypadku twierdzenia Cherna--Gaussa--Bonneta
z rozdziału~\ref{ch:chern-weil}.
\end{proof}

Dla sfery $\chi(S^2)=2$, więc pole maleje z prędkością
$8\pi$, dokładnie jak w Przykładzie~\ref{prz:ricci-sphere}.
Dla torusa $\chi(T^2)=0$ pole jest stałe; metryka
płaska z Przykładu~\ref{prz:ricci-torus} jest szczególnym
rozwiązaniem stacjonarnym. Dla zwartej orientowalnej
powierzchni rodzaju $q>1$ mamy $\chi=2-2q<0$,
więc pole rośnie. Ostatni wniosek dotyczy pola każdego
istniejącego rozwiązania i nie stwierdza jego zbieżności.

\begin{twierdzenie}[Prosta zasada maksimum dla skalara]
\label{thm:ricci-scalar-max-principle}
Niech $g(t)$ będzie gładkim rozwiązaniem
\eqref{eq:ricci-flow} na zwartej rozmaitości bez brzegu,
na przedziale $[0,T)$. Wówczas funkcja
\[
 m(t)=\min_{x\in M}\operatorname{Scal}_{g(t)}(x)
\]
jest niemalejąca. W szczególności początkowa
nieujemność krzywizny skalarnej jest zachowana,
a jeśli $\operatorname{Scal}_{g_0}>0$ w każdym punkcie,
to pozostaje dodatnia w każdej chwili istnienia.
\end{twierdzenie}
\begin{proof}
Ustalmy $0\leq t_0<t_1<T$ i liczby $\varepsilon,\eta>0$.
Na $M\times[t_0,t_1]$ połóżmy
\[
 u(x,t)=\operatorname{Scal}_{g(t)}(x)
       -m(t_0)+\eta+\varepsilon(t-t_0).
\]
W chwili $t_0$ jest $u\geq\eta>0$.
Gdyby później $u$ osiągnęło zero, w pierwszej takiej
chwili $t_*>t_0$ istniałby punkt minimum $x_*$,
w którym $\partial_tu(x_*,t_*)\leq0$ i
$\Delta_{g(t_*)}u(x_*,t_*)\geq0$.
Jednak~\eqref{eq:ricci-scalar-evolution} daje tam
\[
 (\partial_t-\Delta_g)u
   =2|\operatorname{Ric}|^2+\varepsilon>0,
\]
sprzeczność. Zatem $u>0$ na całym cylindrze.
Przechodząc do granicy $\eta,\varepsilon\downarrow0$,
otrzymujemy
$\operatorname{Scal}_{g(t_1)}(x)\geq m(t_0)$
dla każdego $x$. Zwarta przestrzeń zapewnia
istnienie minimum; brak brzegu usuwa potrzebę
warunków brzegowych.
\end{proof}

\section{Poprawka DeTurcka i krótki czas istnienia}
\label{sec:ricci-deturck}
\index{poprawka DeTurcka}
\index{potok Ricciego!krótki czas istnienia}

Równanie~\eqref{eq:ricci-flow} przypomina równanie
ciepła, bo Ricci zawiera drugie pochodne metryki.
Jego główny symbol ma jednak kierunki zerowe
pochodzące od swobody zmiany współrzędnych.
Pokażemy, jak dodatkowa pochodna Liego usuwa tę
degenerację. W tej sekcji istnienie rozwiązania
układu parabolicznego pozostaje jawnym wejściem
analitycznym.

Ustalmy pomocniczą gładką metrykę $\bar g$ na $M$,
na przykład $\bar g=g_0$, i oznaczmy jej symbole
Christoffela przez $\bar\Gamma^k_{pq}$. Różnica
koneksji jest tensorem, więc wzór
\begin{equation}\label{eq:ricci-deturck-vector}
 W^k(g)=g^{pq}\bigl(\Gamma^k_{pq}(g)
                            -\bar\Gamma^k_{pq}\bigr)
\end{equation}
określa globalne pole wektorowe. Składnik
$\bar\Gamma$ wskazuje punkt odniesienia dla
współrzędnych; nie zmienia krzywizny metryki $g$.
\emph{Równanie Ricciego--DeTurcka} ma postać
\begin{equation}\label{eq:ricci-deturck}
 \partial_t\widehat g
  =-2\operatorname{Ric}(\widehat g)
    +\Lie_{W(\widehat g)}\widehat g,
 \qquad \widehat g(0)=g_0.
\end{equation}

Wyprowadźmy część z najwyższymi pochodnymi
w lokalnej mapie. Z definicji symboli Christoffela
i krzywizny, pomijając tylko składniki zawierające
co najwyżej pierwsze pochodne $g$, otrzymujemy
\begin{equation}\label{eq:ricci-principal}
 (-2\operatorname{Ric}(g))_{ij}
 =g^{pq}\bigl(
   \partial_p\partial_q g_{ij}
  +\partial_i\partial_j g_{pq}
  -\partial_p\partial_i g_{qj}
  -\partial_p\partial_j g_{qi}\bigr)
  +\text{wyrazy niższego rzędu}.
\end{equation}
Z kolei po opuszczeniu indeksu pola~\eqref{eq:ricci-deturck-vector}
mamy
\[
 W_j=g^{pq}\bigl(\partial_p g_{qj}
                  -\tfrac12\partial_j g_{pq}\bigr)
      -g_{jk}g^{pq}\bar\Gamma^k_{pq}.
\]
Ponieważ $(\Lie_Wg)_{ij}=\nabla_iW_j+\nabla_jW_i$,
jego część drugiego rzędu to
\begin{equation}\label{eq:ricci-deturck-principal-correction}
 (\Lie_Wg)_{ij}
 =g^{pq}\bigl(
   \partial_i\partial_p g_{qj}
  +\partial_j\partial_p g_{qi}
  -\partial_i\partial_j g_{pq}\bigr)
  +\text{wyrazy niższego rzędu}.
\end{equation}
Pochodne ustalonej $\bar\Gamma$ i pochodne współczynników
$g^{pq}$ dają tylko wyrazy niższego rzędu względem
nieznanej metryki. Dodanie
\eqref{eq:ricci-principal} i
\eqref{eq:ricci-deturck-principal-correction} usuwa
wszystkie mieszane drugie pochodne. Równanie dla
każdej składowej ma więc część główną
\begin{equation}\label{eq:ricci-deturck-heat-principal}
 \partial_t\widehat g_{ij}
     =\widehat g^{pq}\partial_p\partial_q
             \widehat g_{ij}
       +\text{wyrazy niższego rzędu}.
\end{equation}
Po zamrożeniu współczynników w punkcie symbol
przestrzenny prawej strony to
$-g^{pq}\xi_p\xi_q\,\mathrm{Id}$ na tensorach
symetrycznych. Dodatnia określoność $g$ oznacza
$g^{pq}\xi_p\xi_q>0$ dla $\xi\ne0$: układ
\eqref{eq:ricci-deturck} jest ściśle paraboliczny.

\begin{twierdzenie}[Krótki czas istnienia; wejście analityczne]
\label{thm:ricci-short-time}
Dla każdej gładkiej metryki $g_0$ na gładkiej,
zwartej rozmaitości bez brzegu istnieje $T>0$
i dokładnie jedno gładkie rozwiązanie potoku Ricciego
z tą wartością początkową na dostatecznie krótkim
przedziale $[0,T)$. Jednoznaczność dotyczy gładkich
rozwiązań o tej samej wartości początkowej.
\end{twierdzenie}

\begin{proof}[Redukcja do twierdzenia parabolicznego]
\emph{Wejście analityczne} brzmi następująco: gładki
ściśle paraboliczny układ quasiliniowy drugiego rzędu
na zwartej rozmaitości bez brzegu, z gładką wartością
początkową, ma lokalnie w czasie jednoznaczne gładkie
rozwiązanie; ta sama lokalna teoria dotyczy równania
ciepła dla odwzorowań między zwartymi rozmaitościami,
również przy gładkiej metryce dziedziny zależnej od czasu.
Dowód tego twierdzenia wymaga oszacowań parabolicznych
i nie jest zawarty w niniejszym skrypcie. Krótki czas
istnienia potoku udowodnił
\href{https://projecteuclid.org/journals/journal-of-differential-geometry/volume-17/issue-2/Three-manifolds-with-positive-Ricci-curvature/10.4310/jdg/1214436922.full}
{Hamilton, \emph{Three-manifolds with positive Ricci
curvature} (1982), \S\,4}; redukcję do układu ściśle
parabolicznego i równanie cechowania podaje
\href{https://projecteuclid.org/journals/journal-of-differential-geometry/volume-18/issue-1/Deforming-metrics-in-the-direction-of-their-Ricci-tensors/10.4310/jdg/1214509286.full}
{DeTurck, \emph{Deforming metrics in the direction
of their Ricci tensors} (1983)}. Dzięki
\eqref{eq:ricci-deturck-heat-principal} daje ono
rozwiązanie $\widehat g$ równania~\eqref{eq:ricci-deturck}.
Przez ciągłość $\widehat g(t)$ pozostaje dodatnio
określona dla dostatecznie małego $t$.

Niech $\phi_t$ rozwiązuje równanie zwyczajne
\[
 \partial_t\phi_t=-W(\widehat g(t))\circ\phi_t,
 \qquad \phi_0=\id_M.
\]
Gładkość pola i zwartość $M$ dają rodzinę
dyfeomorfizmów w krótkim czasie. Dla
$g(t)=\phi_t^*\widehat g(t)$ wzór na pochodną
cofnięcia tensora daje
\[
 \partial_tg
 =\phi_t^*\bigl(\partial_t\widehat g
                  -\Lie_{W(\widehat g)}\widehat g\bigr)
 =-2\phi_t^*\operatorname{Ric}(\widehat g)
 =-2\operatorname{Ric}(g).
\]
Ostatnia równość wynika z naturalności koneksji
Leviego-Civity i Ricciego względem dyfeomorfizmów.
To konstruuje rozwiązanie oryginalnego równania.

Dla jednoznaczności trzeba jeszcze przenieść
\emph{dowolne} rozwiązanie oryginalne do tego samego
układu parabolicznego. Dla gładkiego potoku $g(t)$
rozwiązujemy równanie ciepła dla odwzorowań
\[
 \partial_tF_t=\tau_{g(t),\bar g}(F_t),
 \qquad F_0=\id_M,
\]
gdzie $\tau$ jest śladem drugiej pochodnej
odwzorowania z $(M,g(t))$ do $(M,\bar g)$.
W mapach jego główna część to
$g^{ij}\partial_i\partial_jF^k$; lokalne istnienie
i jednoznaczność jest drugą częścią wskazanego
wejścia analitycznego. Przez krótki czas $F_t$ jest
dyfeomorfizmem, gdyż $F_0=\id_M$ i $M$ jest zwarte.
Metryka $\widehat g=(F_t^{-1})^*g$ spełnia
\eqref{eq:ricci-deturck}: po przeniesieniu przez
izometrię $F_t:(M,g)\to(M,\widehat g)$ pole
$\partial_tF_t\circ F_t^{-1}$ jest równe
$-W(\widehat g)$, gdyż napięcie identyczności
$(M,\widehat g)\to(M,\bar g)$ ma współrzędne
$\widehat g^{pq}(\bar\Gamma^k_{pq}
-\widehat\Gamma^k_{pq})$.
Dokładniej, dla $V_t=\partial_tF_t\circ F_t^{-1}$
wzór na pochodną cofnięcia daje
\[
 \partial_tg=F_t^*(\partial_t\widehat g
                       +\Lie_{V_t}\widehat g),
 \qquad V_t=-W(\widehat g).
\]
Stąd pochodna przeniesionej metryki zawiera
$+\Lie_{W(\widehat g)}\widehat g$.
Jednoznaczność układu parabolicznego wyznacza
$\widehat g$. Wówczas $F_t$ spełnia zwyczajne
równanie $\partial_tF_t=-W(\widehat g)\circ F_t$
z warunkiem $F_0=\id_M$, więc też jest wyznaczone
jednoznacznie; stąd $g=F_t^*\widehat g$.
Ten krok również
korzysta z podanego twierdzenia analitycznego;
rachunek symbolu sam w sobie nie dowodzi istnienia
ani jednoznaczności PDE.
\end{proof}

\section{Normalizacja objętości i zmiana zegara}
\label{sec:ricci-normalized}
\index{potok Ricciego!znormalizowany}

Potok okrągłej sfery z Przykładu~\ref{prz:ricci-sphere}
zmniejsza jej rozmiar, choć sfera w każdej chwili pozostaje
okrągła. Aby oddzielić zmianę \emph{kształtu} od zmiany
całkowitej objętości, pomnożymy metrykę przez czynnik zależny
od czasu, jednakowy w każdym punkcie. Takie mnożenie wymaga
równoczesnej zmiany parametru czasu.

Niech $V(g)=\int_M\dd\mu_g$ i niech
\begin{equation}\label{eq:ricci-average-scalar}
 r(g)=\frac{1}{V(g)}\int_M\operatorname{Scal}_g\,\dd\mu_g
\end{equation}
oznacza średnią krzywiznę skalarną. Liczba $r(g)$ zależy
od metryki, lecz nie od punktu $M$. Dla $a>0$ ze
Stwierdzenia~\ref{prop:ricci-scaling} wynika
\begin{equation}\label{eq:ricci-average-scaling}
 V(ag)=a^{n/2}V(g),\qquad r(ag)=a^{-1}r(g).
\end{equation}
Rzeczywiście, licznik w~\eqref{eq:ricci-average-scalar}
otrzymuje czynnik $a^{n/2-1}$, a mianownik $a^{n/2}$.

\begin{definicja}[Potok Ricciego znormalizowany objętościowo]
\label{def:ricci-normalized}
Gładką rodzinę metryk $\widetilde g(s)$ nazywamy
\emph{potokiem Ricciego znormalizowanym objętościowo},
jeżeli
\begin{equation}\label{eq:ricci-normalized}
 \partial_s\widetilde g
 =-2\operatorname{Ric}(\widetilde g)
   +\frac{2}{n}r(\widetilde g)\widetilde g.
\end{equation}
Wyraz z $r$ jest wielokrotnością metryki stałą na $M$
w ustalonej chwili; służy zachowaniu całkowitej objętości.
\end{definicja}

\begin{stwierdzenie}[Zachowanie objętości]
\label{prop:ricci-normalized-volume}
Każde rozwiązanie~\eqref{eq:ricci-normalized} spełnia
$V(\widetilde g(s))=V(\widetilde g(0))$. Punktowo zachodzi
\begin{equation}\label{eq:ricci-normalized-volume}
 \partial_s\dd\mu_{\widetilde g}
 =\bigl(-\operatorname{Scal}_{\widetilde g}
          +r(\widetilde g)\bigr)\dd\mu_{\widetilde g}.
\end{equation}
\end{stwierdzenie}
\begin{proof}
Dla dowolnej zmiany $h=\partial_sg$ wzór wyprowadzony
w dowodzie Stwierdzenia~\ref{prop:ricci-volume} ma postać
$\partial_s\dd\mu_g=\frac12\operatorname{tr}_g(h)\dd\mu_g$.
Ślad prawej strony~\eqref{eq:ricci-normalized} wynosi
$-2\operatorname{Scal}+2r$, ponieważ
$\operatorname{tr}_g g=n$. To daje wzór punktowy. Po
scałkowaniu po $M$ otrzymujemy
\[
 \frac{\dd V}{\dd s}
 =-\int_M\operatorname{Scal}\,\dd\mu_g+r(g)V(g)=0
\]
z definicji średniej.
\end{proof}

Zachowanie całki nie oznacza zachowania gęstości objętości
w każdym punkcie: w~\eqref{eq:ricci-normalized-volume}
przyrost zależy od różnicy między lokalnym skalarem i jego
średnią. Normalizacja może więc przenosić objętość między
różnymi obszarami.

\begin{twierdzenie}[Związek obu potoków]
\label{thm:ricci-normalization-clock}
Niech $g(t)$ rozwiązuje~\eqref{eq:ricci-flow} na przedziale
$[0,T)$, a $V_0=V(g(0))$. Połóżmy
\begin{equation}\label{eq:ricci-normalization-clock}
 c(t)=\left(\frac{V_0}{V(g(t))}\right)^{2/n},
 \qquad s(t)=\int_0^t c(u)\,\dd u,
 \qquad \widetilde g(s(t))=c(t)g(t).
\end{equation}
Wtedy $s(t)$ jest ściśle rosnące, $\widetilde g$ spełnia
\eqref{eq:ricci-normalized} i ma objętość $V_0$.
Odwrotnie, każdy potok znormalizowany powstaje lokalnie
w czasie z potoku zwykłego przez taką zmianę skali i czasu.
\end{twierdzenie}
\begin{proof}
Wzór~\eqref{eq:ricci-volume} po scałkowaniu daje
$V'(t)=-r(g(t))V(t)$. Logarytmiczne różniczkowanie
pierwszej formuły~\eqref{eq:ricci-normalization-clock}
prowadzi do
\begin{equation}\label{eq:ricci-normalization-derivative}
 \frac{c'}c=\frac{2}{n}r(g(t)).
\end{equation}
Z~\eqref{eq:ricci-average-scaling} mamy
$r(\widetilde g)=c^{-1}r(g)$. Ricci jako tensor typu
$(0,2)$ nie zmienia się przy mnożeniu metryki przez
stały na $M$ czynnik $c(t)$. Zatem, używając
$\dd s/\dd t=c(t)>0$, liczymy
\[
 \begin{split}
 \frac{\partial\widetilde g}{\partial s}
 &=\frac1c\frac{\dd}{\dd t}(cg)
   =\frac{c'}{c^2}\widetilde g
      -2\operatorname{Ric}(\widetilde g)\\
 &=-2\operatorname{Ric}(\widetilde g)
    +\frac{2}{n}r(\widetilde g)\widetilde g.
 \end{split}
\]
Objętość $cg$ wynosi $c^{n/2}V(g)=V_0$.

Dla odwrotnej konstrukcji wybieramy rozwiązanie
$a'(s)/a(s)=-2r(\widetilde g(s))/n$ z $a(0)=1$,
a następnie definiujemy $t(s)=\int_0^s a(u)\,\dd u$
oraz $g(t(s))=a(s)\widetilde g(s)$. Ponieważ
$\operatorname{Ric}(a\widetilde g)
=\operatorname{Ric}(\widetilde g)$, zróżniczkowanie
ostatniej równości i użycie $t'=a$ usuwa składnik
normalizujący. Otrzymujemy $\partial_tg=-2\operatorname{Ric}(g)$.
\end{proof}

\begin{stwierdzenie}[Ewolucja skalara po normalizacji]
\label{prop:ricci-normalized-scalar}
W potoku~\eqref{eq:ricci-normalized} zachodzi
\begin{equation}\label{eq:ricci-normalized-scalar}
 (\partial_s-\Delta_{\widetilde g})
       \operatorname{Scal}_{\widetilde g}
 =2|\operatorname{Ric}_{\widetilde g}|^2
   -\frac{2}{n}r(\widetilde g)
                   \operatorname{Scal}_{\widetilde g}.
\end{equation}
\end{stwierdzenie}
\begin{proof}
Podstawiamy $h=-2\operatorname{Ric}+(2r/n)g$
do wzoru wariacyjnego~\eqref{eq:ricci-scalar-variation}.
Część $-2\operatorname{Ric}$ daje
$\Delta\operatorname{Scal}+2|\operatorname{Ric}|^2$,
jak w dowodzie Twierdzenia~\ref{thm:ricci-scalar-evolution}.
Ponieważ $r$ jest stałe w przestrzeni, pochodne kowariantne
$(2r/n)g$ znikają. Pozostaje składnik
$-\langle(2r/n)g,\operatorname{Ric}\rangle
=-(2r/n)\operatorname{Scal}$.
\end{proof}

\begin{przyklad}[Metryki Einsteina i powierzchnie]
\label{prz:ricci-normalized-einstein}
Jeśli $\operatorname{Ric}(g_0)=\lambda g_0$, to
$\operatorname{Scal}_{g_0}=n\lambda$ jest stała, więc
$r(g_0)=n\lambda$. Prawa strona
\eqref{eq:ricci-normalized} zeruje się: metryka $g_0$
jest rozwiązaniem stacjonarnym, niezależnie od znaku
$\lambda$. Okrągła sfera zatem nie kurczy się po
normalizacji, choć w zwykłym potoku kurczy się
według~\eqref{eq:ricci-sphere}.

Na zamkniętej powierzchni $\operatorname{Ric}=Kg$,
$\operatorname{Scal}=2K$ i z Gaussa--Bonneta
$\int_M K\,\dd\mu=2\pi\chi(M)$. Ponieważ
$A=\operatorname{Area}(M)$ jest stałe w potoku
znormalizowanym,
\[
 \overline K=\frac{2\pi\chi(M)}{A},\qquad
 \partial_sg=-2(K-\overline K)g.
\]
Jeżeli $K$ w jednym miejscu przewyższa średnią, metryka
tam chwilowo maleje; jeśli jest mniejsze, rośnie.
To jest lokalna interpretacja~\eqref{eq:ricci-normalized-volume}.
Zbieżność potoku powierzchni wymaga dalszych oszacowań.
\end{przyklad}

\section{Iloczyny i lewostronnie niezmiennicza metryka na $S^3$}
\label{sec:ricci-products-homogeneous}
\index{potok Ricciego!metryka iloczynowa}
\index{potok Ricciego!metryka jednorodna}

Obliczenia dla sfery i torusa obejmowały tylko metryki
Einsteina. Iloczyn dwóch czynników może zmieniać ich skale
w różnym tempie; lewostronnie niezmiennicza metryka na
grupie $SU(2)\cong S^3$ pozwoli dodatkowo śledzić zmianę
\emph{kształtu} nawet przy stałym typie topologicznym.

\begin{stwierdzenie}[Potok iloczynu metryk Einsteina]
\label{prop:ricci-product}
Niech $(M_i,h_i)$ będą zwarte, bez brzegu,
$\dim M_i=n_i\geq1$ i
$\operatorname{Ric}(h_i)=\lambda_i h_i$ dla $i=1,2$.
Dla metryki iloczynowej
$g(t)=a(t)h_1\oplus b(t)h_2$ o $a(0)=a_0>0$,
$b(0)=b_0>0$ potok zwykły ma postać
\begin{equation}\label{eq:ricci-product}
 a(t)=a_0-2\lambda_1t,\qquad
 b(t)=b_0-2\lambda_2t
\end{equation}
tak długo, jak oba współczynniki są dodatnie.
Krzywizna skalarna jest stała na $M_1\times M_2$ i wynosi
\begin{equation}\label{eq:ricci-product-scalar}
 \operatorname{Scal}(g(t))
 =\frac{n_1\lambda_1}{a(t)}
  +\frac{n_2\lambda_2}{b(t)}.
\end{equation}
\end{stwierdzenie}
\begin{proof}
Pola z różnych czynników komutują, a metryka iloczynowa
czyni je prostopadłymi. Wzór Koszula pokazuje, że
koneksja Leviego-Civity zachowuje podział styczny
$TM=TM_1\oplus TM_2$: na polach z jednego czynnika
jest koneksją tego czynnika, a pochodne mieszane są zerowe.
W tym rachunku stałe przestrzennie czynniki $a(t),b(t)$
nie zmieniają koneksji. Tensor Ricciego jest więc
\[
 \operatorname{Ric}(g(t))
   =\lambda_1h_1\oplus\lambda_2h_2.
\]
Porównując składowe równania $\partial_tg=-2\operatorname{Ric}$,
dostajemy $a'=-2\lambda_1$ i $b'=-2\lambda_2$.
Kontrakcja Ricciego macierzą odwrotną
$a^{-1}h_1^{-1}\oplus b^{-1}h_2^{-1}$ daje wzór
na skalar. Jednoznaczność z Twierdzenia~\ref{thm:ricci-short-time}
pokazuje, że otrzymana rodzina jest rozwiązaniem potoku
o wskazanej wartości początkowej.
\end{proof}

\begin{przyklad}[Dwie okrągłe sfery]
\label{prz:ricci-two-spheres}
Niech $p,q\geq2$, a $h_p,h_q$ będą metrykami jednostkowych
sfer $S^p,S^q$. Dla
$g_0=\alpha_0h_p\oplus\beta_0h_q$
zapis~\eqref{eq:ricci-product} przyjmuje postać
\begin{equation}\label{eq:ricci-two-spheres}
 \alpha(t)=\alpha_0-2(p-1)t,\qquad
 \beta(t)=\beta_0-2(q-1)t.
\end{equation}
Promienie czynników to odpowiednio $\sqrt\alpha$ i
$\sqrt\beta$. Pierwszy zanika przy
$T_p=\alpha_0/[2(p-1)]$, drugi przy
$T_q=\beta_0/[2(q-1)]$; rozwiązanie istnieje w tej
postaci do $\min(T_p,T_q)$. Jeżeli $T_p<T_q$, pierwszy
czynnik traci skalę, gdy drugi ma jeszcze dodatni promień.
Metryka początkowa iloczynu jest Einsteina dokładnie wtedy,
gdy
\begin{equation}\label{eq:ricci-product-einstein}
 \frac{p-1}{\alpha_0}=\frac{q-1}{\beta_0}.
\end{equation}
Wówczas obie skale zanikają jednocześnie, a potok
znormalizowany jest stacjonarny.
\end{przyklad}

\begin{stwierdzenie}[Równania iloczynu po normalizacji]
\label{prop:ricci-product-normalized}
Przy założeniach Stwierdzenia~\ref{prop:ricci-product},
jeżeli $\widetilde g(s)=a(s)h_1\oplus b(s)h_2$
rozwiązuje potok znormalizowany, to
\begin{equation}\label{eq:ricci-product-normalized}
 \begin{aligned}
 a'&=-2\lambda_1+\frac{2r}{n}a,\\
 b'&=-2\lambda_2+\frac{2r}{n}b,\\
 r&=\frac{n_1\lambda_1}{a}
    +\frac{n_2\lambda_2}{b},\qquad n=n_1+n_2.
 \end{aligned}
\end{equation}
Współczynnik objętości $a^{n_1/2}b^{n_2/2}$ jest stały.
Metryka iloczynowa jest punktem stacjonarnym wtedy
i tylko wtedy, gdy $\lambda_1/a=\lambda_2/b$.
\end{stwierdzenie}
\begin{proof}
Ricci iloczynu ma postać obliczoną wyżej, a jego
skalar jest przestrzennie stały. Wstawienie obu do
\eqref{eq:ricci-normalized} i porównanie składowych
daje układ~\eqref{eq:ricci-product-normalized}.
Ponadto
\[
 \frac{\dd}{\dd s}\log
     \bigl(a^{n_1/2}b^{n_2/2}\bigr)
 =\frac{n_1a'}{2a}+\frac{n_2b'}{2b}
 =-r+\frac{n_1+n_2}{n}r=0.
\]
Oba prawe boki układu są zerem dokładnie wtedy,
gdy $\lambda_1/a=\lambda_2/b=r/n$.
\end{proof}

\begin{przyklad}[Sfera razy okrąg: osobliwość przy stałej objętości]
\label{prz:ricci-sphere-circle-normalized}
Na $S^2\times S^1$ weźmy
$\widetilde g(s)=\alpha(s)h_{S^2}\oplus\beta(s)\dd\theta^2$,
gdzie $h_{S^2}$ ma krzywiznę jeden, a $\theta$ jest
kątem na okręgu. Wtedy $n=3$, $\lambda_1=1$,
$\lambda_2=0$ oraz $r=2/\alpha$. Układ
\eqref{eq:ricci-product-normalized} upraszcza się do
\[
 \alpha'=-\frac23,\qquad
 \frac{\beta'}{\beta}=\frac{4}{3\alpha}.
\]
Po rozwiązaniu z dodatnimi danymi $\alpha_0,\beta_0$
otrzymujemy
\begin{equation}\label{eq:ricci-sphere-circle-normalized}
 \alpha(s)=\alpha_0-\frac23s,\qquad
 \beta(s)=\beta_0
     \left(\frac{\alpha_0}{\alpha(s)}\right)^2,
 \quad 0\leq s<\frac{3\alpha_0}{2}.
\end{equation}
Rzeczywiście,
$\frac{\dd}{\dd s}\log\beta
=-2\alpha'/\alpha=4/(3\alpha)$.
Objętość jest proporcjonalna do
$\alpha\sqrt\beta=\alpha_0\sqrt{\beta_0}$,
lecz krzywizna płaszczyzny stycznej do sfery wynosi
$1/\alpha(s)$ i rośnie bez ograniczenia. Zachowanie
objętości nie wyklucza więc skończonej osobliwości.
Zwarty model należy odróżnić od niezwartych walców
używanych dalej w Przykładzie~\ref{prz:poincare3-walec}.
\end{przyklad}

\subsection*{Metryka Bergera na $SU(2)$}
\label{subsec:ricci-berger}
\index{metryka Bergera}

Niech $X_1,X_2,X_3$ będzie bazą lewostronnie
niezmienniczych pól na $SU(2)$, znormalizowaną przez
\[
 [X_1,X_2]=2X_3,\qquad
 [X_2,X_3]=2X_1,\qquad
 [X_3,X_1]=2X_2.
\]
Jeśli $\omega_1,\omega_2,\omega_3$ jest bazą dualnych
form, to dla $A,C>0$ metryka
\begin{equation}\label{eq:ricci-berger-metric}
 g=A(\omega_1^2+\omega_2^2)+C\omega_3^2
\end{equation}
jest jednorodna: lewe przesunięcia grupy zachowują ją.
Przy $A=C$ jest okrągła; przy $A\ne C$ długości
w kierunku $X_3$ i w dwóch pozostałych kierunkach
są różne. Otrzymujemy przykład, w którym sama
jednorodność nie oznacza metryki Einsteina.

\begin{stwierdzenie}[Równania Bergera]
\label{prop:ricci-berger}
Dla potoku zwykłego startującego z~\eqref{eq:ricci-berger-metric}
metryka ma tę samą postać, a współczynniki spełniają
\begin{equation}\label{eq:ricci-berger-ode}
 A'=-8+4\frac CA,\qquad
 C'=-4\frac{C^2}{A^2}.
\end{equation}
Iloraz $\rho=C/A$ spełnia
\begin{equation}\label{eq:ricci-berger-ratio}
 \rho'=\frac{8}{A}\rho(1-\rho).
\end{equation}
W szczególności $\rho$ rośnie, gdy $0<\rho<1$,
a maleje, gdy $\rho>1$, dopóki metryka jest dodatnio
określona.
\end{stwierdzenie}
\begin{proof}
W ustalonej chwili niech
$e_1=X_1/\sqrt A$, $e_2=X_2/\sqrt A$,
$e_3=X_3/\sqrt C$. Jest to rama ortonormalna.
Oznaczmy $u=2/\sqrt C$ i $v=2\sqrt C/A$. Wtedy
\[
 [e_1,e_2]=v e_3,\qquad
 [e_2,e_3]=u e_1,\qquad
 [e_3,e_1]=u e_2.
\]
Wzór Koszula dla pól lewostronnie niezmienniczych,
których iloczyny skalarne są stałe, redukuje się do
\[
 2\langle\nabla_{e_i}e_j,e_k\rangle
 =\langle[e_i,e_j],e_k\rangle
  -\langle[e_j,e_k],e_i\rangle
  +\langle[e_k,e_i],e_j\rangle.
\]
Podstawiając trzy nawiasy, otrzymujemy między innymi
\[
 \nabla_{e_1}e_2=\frac v2e_3,\quad
 \nabla_{e_2}e_3=\frac v2e_1,\quad
 \nabla_{e_3}e_2=\left(-u+\frac v2\right)e_1,
\]
\[
 \nabla_{e_1}e_3=-\frac v2e_2,\qquad
 \nabla_{e_3}e_1=\left(u-\frac v2\right)e_2,
 \qquad \nabla_{e_i}e_i=0.
\]
Na przykład definicja~\eqref{eq:riemann-def} daje
\[
 \begin{split}
 K(e_1,e_2)
 &=\langle R(e_1,e_2)e_2,e_1\rangle\\
 &=0-\frac{v^2}{4}
     -v\left(-u+\frac v2\right)
  =uv-\frac{3v^2}{4}
  =\frac4A-\frac{3C}{A^2}.
 \end{split}
\]
Ten sam rachunek dla pozostałych płaszczyzn daje
\[
 K(e_1,e_3)=K(e_2,e_3)=\frac{v^2}{4}=\frac C{A^2}.
\]
Pozadiagonalne składowe Ricciego znikają. Sumując
krzywizny płaszczyzn zawierających dany kierunek,
dostajemy
\[
 \operatorname{Ric}(e_1,e_1)
 =\operatorname{Ric}(e_2,e_2)
 =\frac4A-\frac{2C}{A^2},\qquad
 \operatorname{Ric}(e_3,e_3)=\frac{2C}{A^2}.
\]
W bazie $X_i$ oznacza to
$\operatorname{Ric}(X_1,X_1)=4-2C/A$
oraz $\operatorname{Ric}(X_3,X_3)=2C^2/A^2$.
Równanie $\partial_tg=-2\operatorname{Ric}$ daje
\eqref{eq:ricci-berger-ode}. Wreszcie
\[
 \rho'=\frac{C'A-CA'}{A^2}
 =\frac{-4\rho^2-\rho(-8+4\rho)}{A}
 =\frac8A\rho(1-\rho).
\]
Dla $A=C$ oba równania dają $A'=C'=-4$,
zgodnie z~\eqref{eq:ricci-sphere} przy $n=3$.
Jednoznaczność z Twierdzenia~\ref{thm:ricci-short-time}
identyfikuje otrzymane rozwiązanie ODE z potokiem
o zadanej metryce początkowej.
\end{proof}

Wynik jest zgodny z ogólnymi wzorami na krzywiznę
lewostronnie niezmienniczych metryk
\href{https://doi.org/10.1016/S0001-8708(76)80002-3}
{Milnora, \emph{Curvatures of left invariant metrics
on Lie groups} (1976)}. Normalizacja objętości
dodaje do obu równań~\eqref{eq:ricci-berger-ode}
ten sam względny składnik $(2r/3)A$ i $(2r/3)C$.
W pochodnej ilorazu te wyrazy się znoszą. Dlatego
monotoniczny kierunek zmiany $\rho$ jest taki sam
w obu wersjach potoku, po uwzględnieniu zmiany czasu.

\section{Ewolucja pełnej krzywizny i tensora Ricciego}
\label{sec:ricci-tensor-evolution}
\index{potok Ricciego!ewolucja tensora krzywizny}

Skalar z~\eqref{eq:ricci-scalar-evolution} jest śladem
Ricciego, więc nie opisuje całego tensora krzywizny.
Do oceny, czy rozwiązanie pozostaje gładkie, potrzebne
są równania dla wszystkich jego składowych. Zapis
wymaga szczególnej uwagi na kolejność argumentów:
\emph{nasz} tensor z Definicji~\ref{def:tensor-krzywizny}
to $\mathcal R(X,Y,Z,W)=g(R(X,Y)Z,W)$.

W stałej lokalnej bazie współrzędnej oznaczmy
\begin{equation}\label{eq:ricci-P-convention}
 P_{ijkl}:=\mathcal R(\partial_i,\partial_j,
                          \partial_l,\partial_k)
          =g(R(\partial_i,\partial_j)\partial_l,\partial_k).
\end{equation}
Zatem $P_{ijij}=K(\partial_i,\partial_j)$, jeżeli
te dwa wektory są ortonormalne, oraz
$\operatorname{Ric}_{jl}=g^{ik}P_{ijkl}$.
W szczególności dla krzywizny sekcyjnej równej
stałej $K$ mamy
\begin{equation}\label{eq:ricci-P-constant-curvature}
 P_{ijkl}=K(g_{ik}g_{jl}-g_{il}g_{jk}).
\end{equation}
Odwrócenie ostatnich dwóch argumentów względem
$\mathcal R$ jest niezbędne przy porównywaniu poniższych
wzorów z literaturą.

Laplasjan na tensorach będzie miał znak cieplny:
\[
 \Delta_gT=\operatorname{tr}_g\nabla^2T
          =g^{pq}\nabla_p\nabla_qT
          =-\nabla^*\nabla T.
\]
Ostatni symbol to laplasjan szorstki z
Rozdziału~\ref{ch:hodge-bochner}; w zapisie
$\nabla_p\nabla_qT$ uwzględniamy również koneksję
na indeksie $q$. Na funkcjach definicja zgadza
się z początkiem niniejszego rozdziału.

Zdefiniujmy pomocniczy tensor kwadratowy
\begin{equation}\label{eq:ricci-B-definition}
 B_{ijkl}=g^{pr}g^{qs}P_{piqj}P_{rksl}.
\end{equation}
Indeksy $p,q,r,s$ są sumowane. Każdy składnik $B$
jest iloczynem dwóch składowych krzywizny po
kontrakcjach metryką odwrotną.

\begin{twierdzenie}[Równanie pełnego tensora krzywizny]
\label{thm:ricci-riemann-evolution}
W zwykłym potoku Ricciego, w stałej bazie
współrzędnej, tensor~\eqref{eq:ricci-P-convention}
spełnia
\begin{equation}\label{eq:ricci-riemann-evolution}
 \begin{aligned}
 (\partial_t-\Delta_g)P_{ijkl}
 ={}&2(B_{ijkl}-B_{ijlk}-B_{iljk}+B_{ikjl})\\
 &-g^{pq}\bigl(
       P_{pjkl}\operatorname{Ric}_{qi}
      +P_{ipkl}\operatorname{Ric}_{qj}\\
 &\hspace{49pt}
      +P_{ijpl}\operatorname{Ric}_{qk}
      +P_{ijkp}\operatorname{Ric}_{ql}\bigr).
 \end{aligned}
\end{equation}
Przy tych samych typach indeksów w potoku
znormalizowanym prawa strona otrzymuje dodatkowo
$+(2r/n)P_{ijkl}$.
\end{twierdzenie}
\begin{proof}
Podajemy rachunek różniczkowy, w tym tożsamość
Bianchiego potrzebną do przestawienia pochodnych.
Dla $h=\partial_tg=-2\operatorname{Ric}$ wzór
\eqref{eq:ricci-connection-variation} daje
\[
 \partial_t\Gamma^m_{jl}
 =-g^{mp}\bigl(\nabla_j\operatorname{Ric}_{lp}
               +\nabla_l\operatorname{Ric}_{jp}
               -\nabla_p\operatorname{Ric}_{jl}\bigr).
\]
Ponieważ $P_{ijkl}=g_{km}R^m{}_{l i j}$,
różniczkowanie obu czynników daje
\[
 \partial_tP_{ijkl}
 =g_{km}\bigl(\nabla_i\partial_t\Gamma^m_{jl}
             -\nabla_j\partial_t\Gamma^m_{il}\bigr)
   -2\operatorname{Ric}_{km}R^m{}_{l i j}.
\]
Wstawiamy poprzedni wzór i stosujemy symetrie $P$.
Jeśli dla zwięzłości położymy
\[
 D_{ijkl}
  =\nabla_i\nabla_k\operatorname{Ric}_{jl}
   -\nabla_i\nabla_l\operatorname{Ric}_{jk}
   -\nabla_j\nabla_k\operatorname{Ric}_{il}
   +\nabla_j\nabla_l\operatorname{Ric}_{ik},
\]
to otrzymujemy
\begin{equation}\label{eq:ricci-P-variation-intermediate}
 \partial_tP_{ijkl}
 =D_{ijkl}
  -g^{pq}\bigl(P_{ijkp}\operatorname{Ric}_{ql}
             +P_{ijpl}\operatorname{Ric}_{qk}\bigr).
\end{equation}
Nie przestawiliśmy tu jeszcze dwóch pochodnych
Ricciego; ich komutator zawiera krzywiznę.

Druga tożsamość Bianchiego, użyta już w dowodzie
Lematu~\ref{lem:ricci-contracted-bianchi}, daje po
jednej kontrakcji
\[
 g^{pq}\nabla_pP_{qjkl}
  =\nabla_k\operatorname{Ric}_{jl}
   -\nabla_l\operatorname{Ric}_{jk}.
\]
Przy zamianie kolejności pochodnych na tensorze
typu $(0,4)$ komutator działa na każdy z czterech
indeksów. Na przykład
\[
 \begin{split}
 [\nabla_p,\nabla_i]P_{qjkl}
 ={}&-R^a{}_{qpi}P_{ajkl}
     -R^a{}_{jpi}P_{qakl}\\
    &-R^a{}_{kpi}P_{qjal}
     -R^a{}_{lpi}P_{qjka}.
 \end{split}
\]
Kontraktujemy tę równość przez $g^{pq}$,
podstawiamy~\eqref{eq:ricci-B-definition}
i korzystamy z pierwszej tożsamości Bianchiego.
Po wykonaniu tego samego rachunku z zamienionymi
$i,j$ otrzymujemy dokładnie
\begin{equation}\label{eq:ricci-P-laplacian-intermediate}
 \begin{aligned}
 \Delta_gP_{ijkl}
 &+2(B_{ijkl}-B_{ijlk}-B_{iljk}+B_{ikjl})\\
 &=D_{ijkl}
   +g^{pq}\bigl(P_{pjkl}\operatorname{Ric}_{qi}
              +P_{ipkl}\operatorname{Ric}_{qj}\bigr).
 \end{aligned}
\end{equation}
Odejmowanie~\eqref{eq:ricci-P-laplacian-intermediate}
od~\eqref{eq:ricci-P-variation-intermediate}
usuwa wszystkie drugie pochodne Ricciego i daje
\eqref{eq:ricci-riemann-evolution}. Rozpisanie
tożsamości komutatorowej w tej konwencji można
porównać także z
\href{https://archive.intlpress.com/site/pub/files/_fulltext/journals/ajm/2006/0010/0002/AJM-2006-0010-0002-a002.pdf}
{Cao--Zhu, \emph{A Complete Proof of the Poincaré and
Geometrization Conjectures} (2006), Proposition 1.3.1};
oryginalny wzór podał
\href{https://doi.org/10.4310/jdg/1214436922}
{Hamilton (1982), twierdzenie 7.1}.

Dla potoku znormalizowanego dodatkowa zmiana
$h=(2r/n)g$ jest w każdej chwili stałą zmianą
skali w przestrzeni. Nie zmienia ona koneksji ani
Ricciego typu $(0,2)$, lecz tensor $P$ typu $(0,4)$
skaluje się tak jak $g$. Stąd dodatkowy wyraz
$(2r/n)P$ w podanym równaniu.
\end{proof}

\begin{stwierdzenie}[Równanie tensora Ricciego]
\label{prop:ricci-tensor-evolution}
Jako tensor typu $(0,2)$ Ricci spełnia
\begin{equation}\label{eq:ricci-tensor-evolution}
 (\partial_t-\Delta_g)\operatorname{Ric}_{ij}
 =2P_{p i q j}\operatorname{Ric}^{pq}
  -2\operatorname{Ric}_i{}^p\operatorname{Ric}_{pj}.
\end{equation}
W potoku znormalizowanym wzór ma tę samą postać,
z $t$ zastąpionym przez $s$ i z metryką
$\widetilde g$ we wszystkich kontrakcjach.
\end{stwierdzenie}
\begin{proof}
Kontraktujemy~\eqref{eq:ricci-riemann-evolution}
zgodnie z $\operatorname{Ric}_{ij}=g^{pq}P_{p i q j}$.
Pochodna $g^{pq}$ z~\eqref{eq:ricci-inverse-variation}
dla $h=-2\operatorname{Ric}$ wnosi
$2\operatorname{Ric}^{pq}P_{p i q j}$.
W składniku z $B$ stosujemy symetrie i pierwszą
tożsamość Bianchiego; po kontrakcji pary
pozostaje $2Q_{ij}$, gdzie
$Q_{ij}=P_{p i q j}\operatorname{Ric}^{pq}$.
Cztery wyrazy liniowe z $P$ i Ricciego dają kolejno
$-Q_{ij}$, $-\operatorname{Ric}_i{}^p\operatorname{Ric}_{pj}$,
$-Q_{ij}$ i $-\operatorname{Ric}_i{}^p\operatorname{Ric}_{pj}$.
Dodając $2Q_{ij}$ z pochodnej metryki odwrotnej,
otrzymujemy~\eqref{eq:ricci-tensor-evolution}.
Szczegółowy rachunek kontrakcji podają
\href{https://archive.intlpress.com/site/pub/files/_fulltext/journals/ajm/2006/0010/0002/AJM-2006-0010-0002-a002.pdf}
{Cao--Zhu (2006), Corollary 1.3.2}.
Przy normalizacji dodany składnik $(2r/n)g$
nie ma przestrzennych pochodnych i nie zmienia
Ricciego typu $(0,2)$; otrzymujemy tę samą
formułę tensorową.
\end{proof}

Możemy skontrolować znaki bez odwołania do
zewnętrznego przykładu. Na okrągłej sferze
$\operatorname{Ric}_{ij}=\lambda(t)g_{ij}(t)$,
ale jako tensor typu $(0,2)$ w stałej bazie
jest równy $(n-1)h_{ij}$, gdzie $g(t)=a(t)h$.
Jego pochodna wynosi więc zero. Wzór
\eqref{eq:ricci-tensor-evolution} daje to samo:
$P_{p i q j}\operatorname{Ric}^{pq}
=\operatorname{Ric}_i{}^p\operatorname{Ric}_{pj}$.
Natomiast przy śledzeniu
$\operatorname{Ric}^i{}_j=g^{ip}\operatorname{Ric}_{pj}$
pojawia się także pochodna metryki odwrotnej.

Kontrakcja~\eqref{eq:ricci-tensor-evolution}
i ponowne użycie
$\partial_tg^{ij}=2\operatorname{Ric}^{ij}$
odtwarza
\[
 (\partial_t-\Delta_g)\operatorname{Scal}
 =2|\operatorname{Ric}|^2.
\]
Istotnie, skurczenie prawej strony równania
dla Ricciego wynosi zero, a cały składnik
$2|\operatorname{Ric}|^2$ pochodzi z pochodnej
$g^{ij}$. W ten sposób równanie tensorowe
zgadza się z bezpośrednim dowodem
Twierdzenia~\ref{thm:ricci-scalar-evolution}.

\subsection*{Operator krzywizny w ruchomej ramie}
\label{subsec:ricci-curvature-operator}

Przy obliczaniu norm wygodniej utrzymywać ramę
ortonormalną mimo zmiany metryki. W każdym punkcie
wybieramy $e_i(t)$ tak, aby
\[
 \partial_t e_i=\operatorname{Ric}^{\sharp}(e_i).
\]
Wtedy
$\partial_tg(e_i,e_j)
=-2\operatorname{Ric}(e_i,e_j)
 +\operatorname{Ric}(e_i,e_j)
 +\operatorname{Ric}(e_i,e_j)=0$.
Rama pozostaje więc ortonormalna.

Niech $\mathscr K$ będzie endomorfizmem
$\Lambda^2TM$ określonym przez
\[
 \langle\mathscr K(e_i\wedge e_j),
                 e_k\wedge e_l\rangle=P_{ijkl}.
\]
Na sferze o krzywiźnie $K$ mamy $\mathscr K=K I$.
Po identyfikacji $\Lambda^2T_pM$ z
$\mathfrak{so}(T_pM)$ oznaczmy przez
$\operatorname{ad}(U\wedge V)=[U,V]$ odwzorowanie
$\Lambda^2\mathfrak{so}\to\mathfrak{so}$ i połóżmy
\[
 \mathscr K^\#
  =\operatorname{ad}\circ(\mathscr K\wedge\mathscr K)
                       \circ\operatorname{ad}^{*}.
\]
Jest to kwadratowe wyrażenie w składowych
$\mathscr K$. W ruchomej ramie cztery pochodne
bazy znoszą cztery wyrazy liniowe
$\operatorname{Ric}*P$ w~\eqref{eq:ricci-riemann-evolution}.
Pozostaje zwarta postać równania Hamiltona:
\begin{equation}\label{eq:ricci-curvature-operator-evolution}
 (\partial_t-\Delta_g)\mathscr K
     =2(\mathscr K^2+\mathscr K^\#).
\end{equation}
Normalizacja dodaje tu
$-(2r/n)\mathscr K$, gdyż krzywizna sekcyjna
i operator $\mathscr K$ skalują się odwrotnie
do metryki. Ta postać i normalizacja
$\mathscr K=I$ dla jednostkowej sfery są podane
również w
\href{https://annals.math.princeton.edu/wp-content/uploads/annals-v167-n3-p10.pdf}
{Böhm--Wilking, \emph{Manifolds with positive curvature
operators are space forms} (2008), s.~1080}.

Dla kontroli współczynnika w~\eqref{eq:ricci-curvature-operator-evolution}
na sferze $\mathscr K=KI$ i
$I^\#=(n-2)I$. Równanie daje
$K'=2(n-1)K^2$, dokładnie tak jak
$K=1/[a_0^2-2(n-1)t]$
z~\eqref{eq:ricci-sphere}.

\begin{stwierdzenie}[Nierówność dla normy krzywizny]
\label{prop:ricci-curvature-inequality}
Istnieje stała $C_n>0$ zależna tylko od wymiaru,
taka że w zwykłym potoku
\begin{equation}\label{eq:ricci-curvature-norm-inequality}
 (\partial_t-\Delta_g)|\operatorname{Rm}|^2
 \leq -2|\nabla\operatorname{Rm}|^2
       +C_n|\operatorname{Rm}|^3.
\end{equation}
Norma dotyczy tensora $P$ typu $(0,4)$ w metryce
$g(t)$; zamiana $P$ na $\mathcal R$ nie zmienia normy.
\end{stwierdzenie}
\begin{proof}
W ruchomej ramie metryka na $\Lambda^2TM$ ma
stałe współczynniki. Reguła iloczynu dla laplasjanu
i~\eqref{eq:ricci-curvature-operator-evolution} dają
\[
 (\partial_t-\Delta_g)|\mathscr K|^2
 =-2|\nabla\mathscr K|^2
   +4\langle\mathscr K,
            \mathscr K^2+\mathscr K^\#\rangle.
\]
W skończenie wymiarowej przestrzeni tensorów
algebraicznych działanie kwadratowe spełnia
$|\mathscr K^2+\mathscr K^\#|
\leq c_n|\mathscr K|^2$ dla pewnej stałej $c_n$.
W ramie ortonormalnej każda para antysymetrycznych
indeksów pojawia się w pełnym tensorze dwa razy,
więc $|\operatorname{Rm}|^2=4|\mathscr K|^2$
i tak samo dla norm pochodnych. Po zmianie stałej
otrzymujemy~\eqref{eq:ricci-curvature-norm-inequality}.
\end{proof}

Równanie~\eqref{eq:ricci-curvature-norm-inequality}
daje oszacowanie krzywizny na krótkim odcinku czasu
przez zasadę maksimum. Nie kontroluje jednak
automatycznie wszystkich pochodnych krzywizny;
te są osobnym wejściem analitycznym w kryterium
przedłużenia poniżej.

\section{Kryterium przedłużenia przed chwilą osobliwą}
\label{sec:ricci-continuation}
\index{potok Ricciego!kryterium przedłużenia}

Skończony koniec przedziału istnienia nie jest jeszcze
opisem kształtu osobliwości. Możemy jednak wykazać,
że pełny tensor krzywizny musi wtedy rosnąć bez
ograniczenia. Zwartość $M$ jest tu istotna:
zapewnia globalne maksimum krzywizny oraz
pozwala stosować oszacowania paraboliczne bez
warunków w nieskończoności.

Nazywamy rozwiązanie $g(t)$ na $[0,T)$
\emph{maksymalnym}, jeżeli nie istnieje gładkie
rozwiązanie tego samego równania z tą samą
wartością początkową na $[0,T+\varepsilon)$
dla żadnego $\varepsilon>0$. Z twierdzenia o
krótkim czasie i jednoznaczności
(Twierdzenie~\ref{thm:ricci-short-time})
wynika istnienie takiego maksymalnego przedziału;
można sklejać zgodne lokalne rozwiązania.

\begin{twierdzenie}[Kryterium przedłużenia Hamiltona]
\label{thm:ricci-continuation}
Niech $M$ będzie zwartą rozmaitością bez brzegu,
a $g(t)$ maksymalnym gładkim rozwiązaniem
zwykłego potoku Ricciego na $[0,T)$.
Jeśli $T<\infty$, to
\begin{equation}\label{eq:ricci-continuation}
 \lim_{t\uparrow T}\max_{x\in M}
       |\operatorname{Rm}_{g(t)}(x)|_{g(t)}
 =+\infty.
\end{equation}
Równoważnie: jednostajne ograniczenie pełnej
krzywizny na $M\times[0,T)$ pozwala przedłużyć
potok przez chwilę $T$.
\end{twierdzenie}
\begin{proof}[Dowód przy jawnym wejściu analitycznym]
Oznaczmy
$K(t)=\max_M|\operatorname{Rm}_{g(t)}|$.
Najpierw załóżmy, że
$K(t)\leq K_0$ dla wszystkich $t<T$.
W każdym punkcie $|\operatorname{Ric}|
\leq c_n|\operatorname{Rm}|$ dla stałej zależnej
tylko od wymiaru, ponieważ Ricci jest kontrakcją
pełnego tensora. Z równania $\partial_tg=-2\operatorname{Ric}$
i nierówności Grönwalla dla kwadratu długości
każdego stałego wektora wynika
\begin{equation}\label{eq:ricci-metric-equivalence}
 e^{-2c_nK_0t}g_0
 \leq g(t)\leq e^{2c_nK_0t}g_0.
\end{equation}
Stąd metryka nie może utracić dodatniej określoności
przed skończonym czasem $T$.

Teraz używamy osobnego \emph{wejścia analitycznego}:
globalne oszacowania pochodnych krzywizny
Hamiltona--Shiego. W zakresie potrzebnym tutaj mówią,
że jeśli na zwartej $M$ rozwiązanie ma
$|\operatorname{Rm}|\leq K_0$ na $[0,T)$, to dla
każdego całkowitego $m\geq0$ i $\delta\in(0,T)$
istnieje skończona stała $C_{m,\delta}$ taka, że
\begin{equation}\label{eq:ricci-shi-input}
 \sup_{M\times[\delta,T)}
   |\nabla^m\operatorname{Rm}|_{g(t)}
   \leq C_{m,\delta}.
\end{equation}
Stałe zależą od wymiaru, $m$, $K_0$ i wybranego
odstępu $\delta$ od chwili początkowej; w obecnej
formie można też dopuścić zależność od $T$
i $g_0$. Dowód~\eqref{eq:ricci-shi-input}
wymaga oszacowań parabolicznych i nie wynika
z samej nierówności
\eqref{eq:ricci-curvature-norm-inequality}.
Źródła to
\href{https://projecteuclid.org/journals/journal-of-differential-geometry/volume-17/issue-2/Three-manifolds-with-positive-Ricci-curvature/10.4310/jdg/1214436922.pdf}
{Hamilton, \emph{Three-manifolds with positive Ricci
curvature} (1982), lemat 14.4 i twierdzenie 14.1}
oraz
\href{https://doi.org/10.4310/jdg/1214443292}
{Shi, \emph{Deforming the metric on complete
Riemannian manifolds} (1989)}.

Wyjaśnijmy, dlaczego takie oszacowania wystarczają.
Wzór~\eqref{eq:ricci-connection-variation} daje
po wstawieniu $h=-2\operatorname{Ric}$
\[
 \partial_t\Gamma^k_{ij}
 =-g^{k\ell}\bigl(
      \nabla_i\operatorname{Ric}_{j\ell}
     +\nabla_j\operatorname{Ric}_{i\ell}
     -\nabla_\ell\operatorname{Ric}_{ij}\bigr).
\]
Różnica koneksji $\nabla^{g(t)}-\nabla^{g_0}$
jest tensorem. Na $[\delta,T)$ jej pochodna
czasowa jest ograniczona przez
\eqref{eq:ricci-shi-input} z $m=1$,
a~\eqref{eq:ricci-metric-equivalence} pozwala
porównać normy. Całkowanie po czasie ogranicza
tę różnicę. Różniczkując ten sam wzór kolejno
w ustalonych mapach i używając oszacowań z
wyższymi $m$, dostajemy indukcyjnie ograniczenia
wszystkich pochodnych metryki względem $g_0$.
Równanie $\partial_tg=-2\operatorname{Ric}$
ogranicza także ich pochodne czasowe. Dlatego
$g(t)$ ma gładką granicę $g_T$ przy
$t\uparrow T$; z~\eqref{eq:ricci-metric-equivalence}
granica jest dodatnio określona.

Uruchamiamy Twierdzenie~\ref{thm:ricci-short-time}
z metryki początkowej $g_T$. Sklejenie z
dotychczasowym rozwiązaniem w chwili $T$
daje gładki potok na dłuższym przedziale.
Zaprzecza to maksymalności. To dowodzi, że
$K(t)$ nie może być jednostajnie ograniczone.

Pozostaje wzmocnić nieograniczoność do
granicy w~\eqref{eq:ricci-continuation}.
Nierówność~\eqref{eq:ricci-curvature-norm-inequality}
i zasada maksimum, zastosowana do
$\max_M|\operatorname{Rm}|^2$, dają dla
pewnej stałej $C_n>0$ nierówność porównawczą
\begin{equation}\label{eq:ricci-curvature-ode-bound}
 K(t)\leq
 \frac{K(t_0)}
      {1-C_nK(t_0)(t-t_0)}
\end{equation}
od $t_0$ do chwili, w której mianownik
pozostaje dodatni. Gdyby istniały chwile
$t_j\uparrow T$ z $K(t_j)\leq L<\infty$,
to dla dużego $j$ mielibyśmy
$C_nL(T-t_j)<1/2$. Wzór
\eqref{eq:ricci-curvature-ode-bound}
dawałby $K(t)\leq2L$ dla $t\in[t_j,T)$.
Na wcześniejszym zwartym odcinku czasu
krzywizna również jest ograniczona. Otrzymalibyśmy
jednostajne ograniczenie na $[0,T)$,
już wykluczone. Stąd granica wynosi $+\infty$.
\end{proof}

\begin{wniosek}[Minimalne tempo wzrostu krzywizny]
\label{cor:ricci-minimal-blowup}
Przy założeniach Twierdzenia~\ref{thm:ricci-continuation}
istnieje stała $C_n>0$ taka, że dla każdego $t<T$
\begin{equation}\label{eq:ricci-minimal-blowup}
 \max_M|\operatorname{Rm}|(t)
 \geq \frac{1}{C_n(T-t)}.
\end{equation}
\end{wniosek}
\begin{proof}
Gdyby dla pewnego $t$ zachodziło
$C_nK(t)(T-t)<1$, to mianownik
w~\eqref{eq:ricci-curvature-ode-bound}
pozostawałby dodatni aż do $T$.
Oszacowanie ograniczałoby $K$ na $[t,T)$,
a gładkość potoku ogranicza ją również na
$[0,t]$. To przeczyłoby
Twierdzeniu~\ref{thm:ricci-continuation}.
\end{proof}

Wzór jest zgodny ze skalowaniem: jeśli $g(t)$
spełnia zwykły potok, to dla stałej $c>0$
rodzina $\widehat g(s)=c\,g(s/c)$ też go spełnia,
a $|\operatorname{Rm}|_{\widehat g}
=c^{-1}|\operatorname{Rm}|_g$.
Na okrągłej sferze, jak i na zwykłym iloczynie
$S^2\times S^1$, krzywizna jest rzędu
$(T-t)^{-1}$. Dolne tempo wzrostu
nie określa jednak geometrii obszaru
o dużej krzywiźnie; ta analiza należy do
następnego rozdziału.

Analiza konkretnych modeli osobliwości, zbieżności
i twierdzenia Perelmana wymagają dalszych narzędzi.

\par\smallskip
{\footnotesize\noindent\emph{Dalsza lektura:}
\href{https://projecteuclid.org/journals/journal-of-differential-geometry/volume-17/issue-2/Three-manifolds-with-positive-Ricci-curvature/10.4310/jdg/1214436922.full}
{R.\ S.\ Hamilton, \emph{Three-manifolds with positive Ricci
curvature}, J. Differential Geom. 17 (1982), 255--306},
zwłaszcza \S\,3--4: pierwotne równanie i jego ewolucja;
\href{https://projecteuclid.org/journals/journal-of-differential-geometry/volume-18/issue-1/Deforming-metrics-in-the-direction-of-their-Ricci-tensors/10.4310/jdg/1214509286.full}
{D.\ M.\ DeTurck, \emph{Deforming metrics in the direction
of their Ricci tensors}, J. Differential Geom. 18 (1983),
157--162}: usunięcie degeneracji i redukcja do teorii
parabolicznej.\par}
